% 数论（综述）
% license CCBY4
% type Wiki

本文根据 CC-BY-SA 协议转载翻译自维基百科\href{https://en.wikipedia.org/wiki/Number_theory}{相关文章}

\begin{figure}[ht]
\centering
\includegraphics[width=6cm]{./figures/ae528936f6668ed9.png}
\caption{} \label{fig_SHuLun_1}
\end{figure}
数论是纯数学的一个分支，主要致力于研究整数以及算术函数。数论学家研究素数，以及由整数构造的数学对象（例如有理数）的性质，或作为整数的推广而定义的对象（例如代数整数）。

整数既可以单独研究，也可以作为某些方程的解来研究（丢番图几何）。数论中的问题常常能够通过研究某些解析对象来理解，例如黎曼ζ函数，它以某种方式编码了整数、素数或其他数论对象的性质（解析数论）。另外，还可以研究实数与有理数之间的关系，例如研究无理数如何用分数来逼近（丢番图逼近）。

数论与几何一样，是数学最古老的分支之一。数论的一个特点是，它处理的往往是非常容易理解却极难解决的命题。例如，费马大定理在最初提出后的 358 年才被证明，而哥德巴赫猜想自 18 世纪提出以来至今仍未解决。德国数学家卡尔·弗里德里希·高斯曾经说过：“数学是科学的女王，而数论是数学的女王。”\(^\text{[1]}\)数论一度被视为纯粹数学的典范，认为它在数学之外没有任何应用，直到 20 世纪 70 年代，人们才发现素数可被用作构建公钥密码算法的基础。
\subsection{历史}

数论是数学的一个分支，研究整数及其性质与关系。[2]整数是一个集合，它在自然数集合${1,2,3,\dots}$的基础上扩展，加入了数0，以及自然数的相反数${-1,-2,-3,\dots}$。数论学家不仅研究素数，还研究由整数构造的数学对象（例如有理数）的性质，或定义为整数推广的对象（例如代数整数）。[3][4]

数论与算术密切相关，一些作者甚至将这两个词视为同义词。[5]然而，今天“算术”一词通常指研究数值运算的学科，并扩展到实数范围。[6]在更具体的意义上，数论仅限于研究整数，关注它们的性质和相互关系。[7]传统上，它也被称为高等算术。[8]到 20 世纪初，“数论”这一术语已被广泛采用。[note 1]其中“number”意为整数，既可以指自然数，也可以指全体整数。[9][10][11]

初等数论研究整数的一些方面，这些方面可以用初等方法来探讨，例如利用初等证明。[12]相比之下，解析数论依赖于复数以及来自分析学和微积分的技巧。[13]代数数论则使用代数结构（如域和环）来研究数的性质及其相互关系。几何数论运用几何学的概念研究数。[14]数论的其他分支还包括：概率数论[15]、组合数论[16]、计算数论[17]，以及应用数论，后者研究数论在科学与技术中的应用。[18]
\subsubsection{起源}
\begin{figure}[ht]
\centering
\includegraphics[width=8cm]{./figures/b3c636666bde843e.png}
\caption{} \label{fig_SHuLun_2}
\end{figure}
在有记录的历史中，数的知识存在于古代的美索不达米亚、埃及、中国和印度文明中【19】。已知最早的算术性质的历史发现是普林普顿 322 号泥板，约公元前 1800 年。它是一块破碎的泥板，包含了一张毕达哥拉斯三元组（即整数$(a,b,c)$，满足$a^{2}+b^{2}=c^{2}$的数对）的列表。这些三元组数量众多且数值很大，不可能通过蛮力枚举得到【20】。

该表的布局表明，它是借助于在现代语言中可写作如下的恒等式构造出来的【21】：
$$
\left(\tfrac{1}{2}\left(x-\tfrac{1}{x}\right)\right)^{2}+1
=\left(\tfrac{1}{2}\left(x+\tfrac{1}{x}\right)\right)^{2},~
$$
这一恒等式隐含在古巴比伦日常的习题中【22】。不过也有人提出，该表可能是作为学校习题的数值例子的来源【23】【注2】。

普林普顿 322 泥板是巴比伦数学中唯一现存的证据，可以被视为与现代所谓数论相对应的成果；然而，当时的一种巴比伦代数学要发达得多【24】。

尽管其他文明可能在最初对希腊数学有所影响【25】，但所有关于这种借鉴的证据都出现得相对较晚【26】【27】，因此很可能希腊的arithmētikḗ（数的理论性或哲学性研究）是一种本土传统【28】。

古希腊人对整除性产生了浓厚兴趣。毕达哥拉斯学派赋予完全数与亲和数神秘的性质。毕达哥拉斯的传统还谈及所谓的多边形数或图形数【29】。

欧几里得在其《几何原本》中专门论述了一些属于初等数论的主题，包括素数与整除【30】。他提出了用于计算两个数的最大公约数的欧几里得算法，并给出了一个证明，表明素数的数量是无限的。

晚期古代最重要的 *arithmētikḗ* 权威是亚历山大的丢番图，他大约生活在公元三世纪。他撰写了《算术》(*Arithmetica*)，这是一部题集，其中的问题无一例外是寻找某个多项式方程组的有理数解，通常形式为：

[
f(x,y) = z^{2}
]

或

[
f(x,y,z) = w^{2}.
]

用现代术语来说，**丢番图方程**是指要求有理数或整数解的多项式方程。

---

要不要我继续帮你翻译“印度数论传统”这一部分？
